% Fichier complété par tools/complete_missing_corrections.py.
% Source : CIN/CIN-01-Parametrage/04_RR/04_RR.tex
% Statut : rédigé (1 question(s)).
\begin{corrigebox}[Corrigé rédigé]
\CorrigeQuestion{1}
On pose $\theta=(\vec i_0,\vec i_1)$ pour le pivot 1/0 et
$\varphi=(\vec i_1,\vec i_2)$ pour le pivot 2/1. Les deux pivots ont pour axe
$\vec k_0$ :
\[\vec\Omega_{1/0}=\dot\theta\vec k_0,\qquad
  \vec\Omega_{2/1}=\dot\varphi\vec k_0,\qquad
  \vec\Omega_{2/0}=(\dot\theta+\dot\varphi)\vec k_0.\]
Avec $\overrightarrow{AB}=R\vec i_1$ et $\overrightarrow{BC}=L\vec i_2$,
$\overrightarrow{AC}=R\vec i_1+L\vec i_2$ et l'orientation absolue de 2 vaut
$\theta+\varphi$.
\end{corrigebox}
